\section{A deterministic calculus with auxiliary formulas}
\label{sec:composition}

The construction combines two restrictions that have different roles.  The
source of an equality elimination must not be the literal term $0$, and the
root of a paired matrix determines its introduction rule.  There is no
independent case/progress choice.  The interaction is resolved by allowing
auxiliary formulas in proofs.  Only the public formulas below are assigned
the ordinary conditional reading in the bi-interpretation theorem.  In
particular, we do not assert a round-trip theorem for every auxiliary
formula.

\subsection{Terms, public formulas, and proof formulas}

There is one numerical sort.  The terms are
\[
 t ::= x\mid 0\mid S(t)\mid I(n,a;u.h),
\]
where $u$ is bound only in $h$.  Substitution is capture-free.  Write
\[
 1:=S0,\qquad 2:=S1,\qquad C(v,a,b):=I(v,a;u.b),
\]
with $u$ fresh in the last abbreviation.  For an equation $E\equiv(s=t)$,
put
\[
 G_v E := \bigl(C(v,0,s)=C(v,0,t)\bigr).
\]

The public language $\G$ is the guarded language of the inherited package:
every equation is public; $\D_{x,y}B$ is public if $B$ is a compound public
formula; and an atomic matrix immediately under $\D_{x,y}$ must have the
literal shape
\begin{equation}
 B(x,y)\equiv I(y,0;u.h)=I(y,0;u.k).
 \tag{G}\label{eq:public-guard}
\end{equation}
The bodies $h,k$ may contain $x,y$, the state variable, and fixed
parameters.  This is a syntactic condition.  Every such matrix has a
provable zero instance.  The designated contradiction is $0=1$.

The auxiliary proof language $\U$ is larger:
\[
 A ::= (s=t)\mid\Q_y A\mid\D_{x,y}A.
\]
Its paired matrices need not satisfy \eqref{eq:public-guard}.  Contexts are
finite lists of reusable auxiliary assumptions.  Public theoremhood means
that a public conclusion has a derivation in this auxiliary calculus; it
does not require every intermediate formula to be public.  The same
definition applies to consequence from public assumptions.

\subsection{All primitive axioms and rules}

The only nonlogical axiom schemas, available under every context, are
\[
 I(0,a;u.h)=a,\qquad
 I(Sn,a;u.h)=h[I(n,a;u.h)/u].
 \tag{I0,IS}
\]
The hypothesis rule is $\Gamma\vdash A$ whenever $A\in\Gamma$.
The remaining rules are
\[
 \frac{\Gamma\vdash s=t\qquad\Gamma\vdash F(s)}
      {\Gamma\vdash F(t)}
 \quad(s\text{ is not literally }0),
 \tag{Tr$^{-}$}\label{eq:restricted-tr}
\]
\[
 \frac{\Gamma\vdash\Q_y B(y)\qquad\Gamma\vdash B(0)}
      {\Gamma\vdash B(t)},
 \tag{QE}
\]
\[
 \frac{\Gamma\vdash\D_{x,y}B(x,y)\qquad
       \Gamma,\Q_yB(a,y)\vdash A}
      {\Gamma\vdash A},
 \tag{DE}
\]
and one root-determined introduction schema
\begin{equation}
 \frac{\Gamma,B(t,\eta_B(z))\vdash B(t,Sz)}
      {\Gamma\vdash\D_{x,y}B(x,y)},\qquad
 \eta_B(z)=
 \begin{cases}
 z,&B\text{ is a literal equation},\\
 0,&B\text{ has root }\Q\text{ or }\D.
 \end{cases}
 \tag{DI$_{\rm root}$}\label{eq:root-di}
\end{equation}
Here $F$ is an auxiliary formula template.  In DE, $a$ is fresh for the
retained context, the conclusion, and the fixed parameters.  In DI,
$z$ is fresh for the retained context, the witness term $t$, and the
fixed parameters.  The displayed local assumption is discharged.
Bound names are always chosen to avoid capture.

There is no primitive equality reflexivity or successor axiom, no
separate induction schema, no unrestricted atomic case introduction,
and no primitive introduction rule for $\Q$.  Weakening, renaming,
contraction, and replacement of a retained assumption by its proof are
admissible by the usual induction on a finite derivation.  No unrestricted
numerical-substitution theorem for auxiliary derivations is needed below:
a substitution could turn an allowed source variable into the literal
$0$, so a direct constructor-wise substitution argument would be invalid.

\subsection{Induction and cases are available before equality is repaired}

\begin{lemma}[Auxiliary equation induction]\label{lem:aux-ei}
For every literal equation $E(y)$, the rule
\[
 \frac{\Gamma\vdash E(0)\qquad
       \Gamma,E(z)\vdash E(Sz)}
      {\Gamma\vdash E(t)}
\]
is admissible, with the usual fresh $z$.  The retained context may contain
arbitrary auxiliary formulas.
\end{lemma}
\begin{proof}
Apply atomic DI to $\D_{x,y}E(y)$, with $x$ unused and witness $0$.
Its premise is precisely the supplied step.  DE exports $\Q_yE(y)$;
the fresh witness is absent from this formula.  QE and the supplied base
give $E(t)$.  This derivation uses no transport at all.  Its paired
intermediate formula is generally auxiliary rather than public.
\end{proof}

\begin{lemma}[Numerical cases]\label{lem:aux-cases}
From $\Gamma\vdash A(0)$ and $\Gamma\vdash A(Sz)$, with $z$ fresh,
one may infer $\Gamma\vdash A(t)$ for every auxiliary $A(y)$.
\end{lemma}
\begin{proof}
For an equation, use Lemma~\ref{lem:aux-ei} and ignore its local induction
assumption.  For a compound formula, case DI with an unused witness
gives $\D_{x,y}A(y)$; DE and QE use the given base to obtain $A(t)$.
\end{proof}

This is the decisive bootstrap.  We do not first try to pad the hypothesis
$0=t$ into a public compound formula.  Auxiliary equation induction is
already available, and supplies the atomic cases needed by the
literal-zero-source argument.

These are the compatibility steps needed to combine the two previously
separate reductions.  The guarded public grammar is inherited from
\cite{guarded}; the source restriction is inherited from \cite{source}.

\subsection{Recovering equality from the restricted source rule}

The following proof reconstructs the zero-source argument in the pinned
\texttt{source\_dispatch} package and
\texttt{NativeZeroSource.lean}.  Its only non-equational ingredients are
equation induction and numerical cases, now supplied by
Lemmas~\ref{lem:aux-ei} and~\ref{lem:aux-cases}.  The formulas in the
calculations below are equations.
The new proof reuses this numerical argument only; the final unrestricted
auxiliary-formula transport argument of the source package is not used
\cite{source,lean-zero}.

\begin{lemma}[Safe equality calculations]\label{lem:safe-equality}
Reflexivity is derivable for every term.  Symmetry is derivable for $p=q$
when $p$ is not literally $0$, and transitivity is derivable through any
middle term that is not literally $0$.  Congruence and replacement on
either side of an equation are available along an equation with such a
source.  If $p,q$ are not literally $0$, then
\[
 \Gamma\vdash p=0,\quad\Gamma\vdash q=0
 \quad\Longrightarrow\quad\Gamma\vdash p=q.
 \tag{Zero join}\label{eq:zero-join}
\]
\end{lemma}
\begin{proof}
Put $r=I(0,t;u.0)$.  The axiom $r=t$, used as both premises of
Tr$^{-}$ with template $v=t$, gives $t=t$.  The usual equality
templates then give symmetry at an allowed source, transitivity through
an allowed middle term, and congruence at an allowed source.

For the last assertion, let $W(v)=I(0,v;u.0)$ and $w=W(0)$.  From
$p=0$ and $W(p)=W(p)$, replace $p$ only on the left to obtain
$w=W(p)$.  Symmetry at $w$, followed by the computation $W(p)=p$
on the left, gives $p=w$.  Similarly $q=w$.  Symmetry at $q$ and
transitivity through $w$ give $p=q$.  All these sources and the middle
term $w$ are syntactically nonzero; no numerical inequality is asserted.
\end{proof}

\begin{lemma}[Zero-source symmetry]\label{lem:zero-symmetry}
For every term $t$ and every auxiliary context,
\[
 \Gamma\vdash 0=t\quad\Longrightarrow\quad\Gamma\vdash t=0.
\]
\end{lemma}
\begin{proof}
The case where $t$ is literally $0$ is reflexivity.  Otherwise $t$ is an
allowed source.  Write
\[
 A(n)=I(n,t;u.Su).
\]
The zero computation and safe symmetry give $t=A(0)$, which may replace
the right side of the assumed $0=t$; thus $0=A(0)$.  Under
$z=A(z)$, congruence at the variable source $z$ gives
$Sz=S(A(z))$, and the successor computation gives $Sz=A(Sz)$.
Auxiliary equation induction therefore gives $1=A(1)$ and hence
\begin{equation}
 1=St.
 \label{eq:one-successor-t}
\end{equation}

For completeness, the small test used to remove this successor is
\[
 \begin{aligned}
 Z(v)&=C(v,1,0),& K(v)&=C(v,0,1),& H(v)&=C(v,1,2),\\
 N(n)&=I(n,0;u.Z(u)),& R(n)&=I(n,0;u.H(u)),& L(n)&=N(R(n)).
 \end{aligned}
\]
Numerical cases and computation give $H(v)=S(K(v))$.  Consequently
$R(1)=1$ and $R(SSv)=2$.  Computation gives $N(1)=1$ and $N(2)=0$,
so $L(1)=1$ and $L(SSv)=0$.  A second numerical case argument yields
\begin{equation}
 L(Sv)=Z(v).
 \label{eq:singleton-test}
\end{equation}
At the zero branch both sides reduce to $1$; at the successor branch
both sides reduce to $0$, and \eqref{eq:zero-join} joins them.  All
reductions use iterator sources.  Similarly numerical cases prove the
reversed filter identity
\begin{equation}
 C(Z(v),v,0)=v.
 \label{eq:reversed-filter}
\end{equation}
Its zero branch reduces via $Z(0)=1$, and its successor branch via
$Z(Sw)=0$.  Again every major used for reduction has a nonzero
syntactic source.

Applying $L$ to \eqref{eq:one-successor-t}, using
$L(1)=1$ and \eqref{eq:singleton-test}, and using safe symmetry and
transitivity, gives $Z(t)=1$.  Hence the left side
$q=C(Z(t),t,0)$ of \eqref{eq:reversed-filter} is also provably $0$.
The filter gives $q=t$.  Symmetry at the iterator source $q$, followed
by transitivity through $q$, yields $t=0$.

Every identity just used is a schema constructed at the displayed
terms.  The argument does not assume that arbitrary restricted
derivations may be numerically substituted.
\end{proof}

\begin{corollary}[Unrestricted equational transport]
Symmetry, transitivity, term congruence, and transport in every equation
template are admissible without a source restriction.
\end{corollary}
\begin{proof}
Symmetry follows from Lemmas~\ref{lem:safe-equality} and
\ref{lem:zero-symmetry}.  For transitivity through literal $0$, the
cases of a literal-zero endpoint are immediate; otherwise reverse the
second equation and use zero join.  For congruence along $0=t$ with
$t$ not literally $0$, reverse to $t=0$, use safe congruence, then
use unrestricted symmetry.  If $F(v)$ is $a(v)=b(v)$, term
congruence, symmetry, and transitivity transform $F(s)$ to $F(t)$.
\end{proof}

\begin{lemma}[Transport in public formulas]\label{lem:public-transport}
The unrestricted transport rule is admissible whenever its template $F$
is public.  Thus numerical equality has its full ordinary use throughout
the public language.
\end{lemma}
\begin{proof}
Induct on the public formula template.  The equation case is the preceding
corollary.  Suppose first that $F(v)=\D_{x,y}B(v,x,y)$ has compound
matrix.  Open $F(s)$ at a fresh witness $a$, obtaining the clause
$\Q_y B(s,a,y)$.  Introduce $F(t)$ at witness $a$ by compound DI.
Its base assumption $B(t,a,0)$ transfers, using the induction hypothesis
and $t=s$, to $B(s,a,0)$.  QE gives $B(s,a,Sz)$; the induction
hypothesis and $s=t$ transfer this to $B(t,a,Sz)$.

If $B$ is atomic, its public guard makes $B(s,a,0)$ a theorem.
The opened source clause therefore gives $B(s,a,Sz)$ outright.
Unrestricted equational transport gives $B(t,a,Sz)$.  Atomic DI
introduces the target, ignoring its local induction assumption.
Close DE in both cases.  Eigenvariables are chosen fresh for all fixed
parameters.
\end{proof}

This lemma is intentionally about public templates.  We do not need, and
do not assert, unrestricted zero-source transport in every auxiliary
formula containing an unguarded atomic paired matrix.

\subsection{Soundness of every auxiliary derivation}

Give auxiliary formulas the following ordinary first-order readings over
the scalar terms:
\[
 \begin{aligned}
 (s=t)^\dagger&=(s=t),\\
 (\Q_y B)^\dagger&=B^\dagger(0)\to\forall y\,B^\dagger(y),\\
 (\D_{x,y}B)^\dagger&=
 \begin{cases}
 \exists x\,\forall y\,(B(x,y)\to B(x,Sy)),&B\text{ is an equation},\\
 \exists x\,(B^\dagger(x,0)\to\forall y\,B^\dagger(x,y)),
       &B\text{ is compound}.
 \end{cases}
 \end{aligned}
\]
Numeric substitution preserves the root category of a formula, so this
reading commutes with capture-free numeric substitution.  The reading of
$\Q$ is always conditional, including at an equation.

\begin{theorem}[Auxiliary soundness]\label{thm:aux-soundness}
If $\Gamma\vdash_{\U}A$, then iPRA proves $A^\dagger$ from
$\Gamma^\dagger$, after the inherited primitive-recursive realization
of scalar terms.
\end{theorem}
\begin{proof}
Induct on the derivation.  Hypothesis, computation, transport, and QE
have their ordinary sound interpretations.  Atomic DI directly proves
the displayed progressive reading.  Compound DI uses numerical coverage
to extend its zero-to-successor premise to a zero-to-all conclusion.
For DE at an atomic matrix, ordinary equation induction converts
$\forall y(E(y)\to E(Sy))$ into $E(0)\to\forall yE(y)$, which is
exactly the exported $\Q$ assumption.  This uses only quantifier-free
induction, with the existential witness and other parameters fixed.
DE at a compound matrix is ordinary existential elimination.
All side assumptions and eigenvariable conditions are preserved.
\end{proof}

Let $A^\circ$ denote the \emph{uniform ordinary reading} on public
formulas: equality is actual numerical equality and every public
$\D_{x,y}B$ means $\exists x(B^\circ(x,0)\to\forall yB^\circ(x,y))$.

\begin{lemma}[Agreement on public formulas]\label{lem:public-agreement}
For every public $A$, iPRA proves $A^\dagger\leftrightarrow A^\circ$.
\end{lemma}
\begin{proof}
At a guarded atomic matrix $E(x,y)$, the zero instance $E(x,0)$ is a
theorem.  Progress implies all instances by equation induction, while
conditional closure implies all instances by that true base and hence
implies progress.  These two readings are therefore equivalent, with
$x$ fixed, and remain equivalent under the existential quantifier.
At compound matrices use structural induction.  Equations are unchanged.
\end{proof}

\subsection{Completeness and both interpretation composites}

We record the inherited guarded translation explicitly, because it
identifies the exact scope of the result.  Let $\mathcal C$ be the
existing native scalar calculus.  Define
\[
 L(E)=\D_{u,v}G_vE,\qquad
 P(A)=\begin{cases}L(A),&A\text{ an equation},\\A,&A\text{ compound}.
 \end{cases}
\]
The new variables in $L$ are fresh.  Translate old Eq/$\D$ formulas by
\[
 T(s=t)=(s=t),\qquad
 T(\D_{x,y}B)=\D_{x,y}P(TB).
\]
If the old presentation also treats $\Q$ as an object constructor, put
$T(\Q_yB)=\D_{x,y}P(TB)$ with $x$ unused.  Let $J$ be the literal
inclusion of public formulas into the old language.

\begin{lemma}[Guarded translation and its composites]\label{lem:guarded-composites}
The translation $T$ preserves derivations of the old calculus in the new
public consequence relation.  Uniformly under additional contexts,
\[
 A\dashv\vdash_{\mathcal C}JT(A),\qquad
 B\dashv\vdash_{\U}TJ(B)\quad(B\text{ public}).
\]
\end{lemma}
\begin{proof}
After equality repair, $L(E)\dashv\vdash E$: from $E$, atomic DI
proves $L(E)$ by ignoring its local assumption; conversely DE, the
theorem $G_0E$, and QE at $1$ recover $E$.  Therefore $P(A)$ is
always compound and interderivable with $A$.

Old case introduction transfers to compound DI after padding the base
and the successor conclusion.  Old equation progress transfers by
auxiliary equation induction under the padded base assumption, followed
by compound DI.  The old direct $\Q$ progress rule is covered by
Lemma~\ref{lem:aux-ei} and the unused-witness translation of $\Q$.
Elimination transfers by replacing each use of its opened clause with
padding of the antecedent, use of the translated clause, and unpadding
of the consequent.  Transport transfers by
Lemma~\ref{lem:public-transport}.  Computation and assumptions are
unchanged.  This is an effective transformation of finite proofs.

The old calculus has binder congruence and the same padding equivalence,
so structural induction gives its round trip.  In the new calculus,
uniform equivalence of two compound matrices lifts under $\D$ by DE
and compound DI.  At a guarded atomic matrix $B$, the remaining step is
\[
 \D_{x,y}B\dashv\vdash\D_{x,y}L(B).
\]
The forward direction opens the source clause, unpads the target base,
uses the clause, and pads the successor.  For the reverse direction,
the true guarded base proves $L(B(a,0))$; the opened clause then gives
$L(B(a,Sz))$ and hence $B(a,Sz)$.  Atomic DI introduces $\D B$,
ignoring its local assumption.  These are all the cases in the new
round-trip induction.  An unused existential witness introduces or
eliminates the old $\Q$ interpretation, using the witness $0$.
\end{proof}

\begin{theorem}[Exact public consequence relation]\label{thm:exact-public}
For public $\Gamma,A$,
\[
 \Gamma\vdash_{\U}A
 \quad\Longleftrightarrow\quad
 \Gamma\vdash_{\mathcal C}A
 \quad\Longleftrightarrow\quad
 \Gamma^\circ\vdash_{\ipra}A^\circ.
\]
\end{theorem}
\begin{proof}
Auxiliary soundness and public agreement give the rightward implication
from the new calculus to iPRA.  The pinned native--iPRA equivalence gives
the old native derivation.  Conversely, translate an old derivation by
$T$ and use both public round-trip equivalences, including the uniform
equivalences for the retained assumptions, to recover the original public
sequent.  The second equivalence is the inherited native--iPRA bridge.
\end{proof}

\begin{corollary}[Bi-interpretability and equiconsistency]
The ordinary realization of this public theory is bi-interpretable and
equiconsistent with full first-order iPRA.  Both numerical domains are the
whole domain; equality is actual equality; and both numerical composites
have the identity graph $x=y$.
\end{corollary}
\begin{proof}
Compose $T,J$ and their formula composites with the pinned
bi-interpretation of the native calculus and iPRA.  No term or numerical
sort is changed by $T$ or $J$, so the numerical composites remain the
identity.  The distinguished contradiction $0=1$ is unchanged and has
the ordinary reading equivalent to falsity in iPRA.  Exact public
consequence therefore gives equiconsistency as well.
\end{proof}

\subsection{Why the auxiliary distinction is essential}

\begin{proposition}[The canonical-only composition would fail]
If paired matrices throughout derivations were required to satisfy the
public grammar, with $\Q$ resources only the corresponding opened
clauses, the simultaneous restrictions would not even prove $1=0$ from
$0=1$.
\end{proposition}
\begin{proof}
Interpret terms in the standard natural numbers.  At the root of a
formula, declare every literal equation $0=t$ true, and interpret every
other equation by actual equality.  Interpret public compounds by their
ordinary conditional reading.  Every atomic leaf under a public paired
binder has an iterator on the left; consequently those leaves always
receive actual equality, and their readings commute with substitution.

Computations, guarded atomic progress, compound case introduction,
and the associated clause elimination are sound for this interpretation.
Restricted transport is sound too.  Its major $s=t$ has a syntactically
nonzero left side, so it asserts actual equality.  If an equation
template instantiated at $s$ has literal $0$ on the left, that left
template was constant $0$: it could not have been the hole, since
$s$ is not literally $0$.  The target is then also true by convention.
Otherwise transport preserves actual equality, unless its target has
literal-zero left side and is true by convention anyway.  Transport
inside a public compound is ordinary extensional transport.

The assumption $0=1$ is true in this interpretation, while $1=0$ is
false.  This proves the stated nonderivability.  The auxiliary induction
used in Lemma~\ref{lem:zero-symmetry} is exactly the extra proof resource
that invalidates this counterinterpretation.
\end{proof}

\begin{remark}[Exact checking contract]
The compact checker is a checker of auxiliary derivations with a
\emph{canonical public-goal contract}.  Its soundness and completeness
claim is Theorem~\ref{thm:exact-public}, for a goal in $\G$.  A raw
auxiliary goal need not be public, even when the program accepts its
proof.  The separately supplied strict checker includes and charges an
explicit public-goal validator.  A claim about the compact kernel alone
must not silently include the validator, and a claim about the whole
auxiliary syntax must not be substituted for the public bi-interpretation
proved here.
\end{remark}
